\chapter{Visual Perception and Self-Healing Execution}
\label{chap:vision_self_healing}
\minitoc
\newpage

\section{Introduction}
\label{sec:ch4_intro}

After presenting the semantic intelligence layer of \textbf{VisionTest-AI}, this chapter focuses on the visual perception layer. While the NLP module is responsible for understanding what the user wants to test, the vision module is responsible for observing the application interface and detecting the elements that can be used during execution.

Traditional web automation tools mainly interact with the Document Object Model (DOM). This approach is powerful when selectors are stable, but it becomes fragile when the interface changes, when elements are dynamically generated, or when the tested page contains components that are difficult to locate using classical selectors. For this reason, \textbf{VisionTest-AI} introduces a visual fallback mechanism based on object detection and OCR.

The objective of this chapter is to present the complete construction of the vision module. It describes the dataset preparation process, including the difficulties encountered during data collection and annotation. It then explains the training of the YOLOv8n model, the OCR pipeline, the vision API implementation, and the self-healing execution mechanism. Finally, it presents realistic observations, successful cases, failed cases, and limitations identified during testing.

\subsection*{Chapter Overview}
\phantomsection
\addcontentsline{toc}{subsection}{Chapter Overview}

This chapter follows the complete visual workflow. It starts with dataset preparation and the annotation difficulties encountered during the project. It then presents YOLOv8n training, OCR preprocessing, the vision API, the self-healing execution mechanism, and the real observations obtained during testing.

\section{Objective of the Vision Module}
\label{sec:vision_objective}

The vision module was designed to make the agent more resilient when classical selector-based automation fails. In a normal execution flow, the system first tries to interact with the page using Playwright selectors generated by the semantic mapping module. However, if these selectors fail, the system needs another way to locate the target element.

The main objective of the vision module is therefore to analyze screenshots and detect visible UI elements such as buttons, input fields, links, dropdowns, and checkboxes. Once these elements are detected, OCR is used to extract visible text and improve the matching between the semantic target and the visual elements.

The module has four main responsibilities:

\begin{itemize}
	\item Detect UI elements from screenshots using YOLOv8n.
	\item Extract visible text using OCR.
	\item Match the semantic target from the NLP module with the detected visual elements.
	\item Provide coordinates that can be used by Playwright for fallback execution.
\end{itemize}

This approach does not replace Playwright selectors. Instead, it complements them. The final execution strategy is based on two plans: Plan A uses classical selectors, while Plan B uses visual detection and coordinate-based interaction when selectors fail.

\section{Dataset Preparation: From Failed Attempts to DOM-Assisted Labeling}
\label{sec:dataset_preparation}

The preparation of the vision dataset was one of the most difficult and time-consuming parts of the project. Unlike common object detection tasks, where public datasets are often available, UI element detection requires screenshots containing realistic web components such as buttons, input fields, links, dropdowns, and checkboxes.

At the beginning of the project, the objective seemed simple: collect enough screenshots and annotate the visible interface elements. However, this phase quickly revealed several practical difficulties. Ready-made datasets were not diverse enough, manual annotation was slow, and automatic labeling tools were not immediately reliable for complex web interfaces.

For this reason, the dataset preparation evolved through several attempts before reaching the final solution based on DOM-assisted auto-labeling.

\subsection{The Challenge of Building a UI Dataset}
\label{subsec:dataset_challenge}

Building a dataset for UI element detection is different from building a dataset for natural objects. In traditional object detection, classes such as cars, persons, or animals have relatively stable visual characteristics. In web interfaces, the same class can appear in many different forms.

For example, a button can be a filled rectangle, an outlined element, a rounded component, a text-only link styled as a button, or an icon button. Similarly, input fields can have different borders, placeholders, icons, labels, and colors depending on the design system used by the website.

This visual diversity makes the dataset preparation phase critical. If the dataset is too narrow, the trained model may only recognize elements similar to those seen during training and fail on unseen websites.

\subsection{First Attempt: Searching for Ready-Made Datasets}
\label{subsec:ready_made_datasets}

The first attempt consisted of searching for existing datasets containing annotated UI elements. The goal was to save time by reusing a public dataset instead of building one from scratch.

However, this approach was not sufficient for the project. Most available datasets were either focused on mobile applications, contained synthetic interfaces, or did not include enough variation in real web layouts. Some datasets also used class names that were not aligned with the needs of \textbf{VisionTest-AI}.

The main limitation was the lack of diversity. A model trained on a limited dataset may perform well on similar screenshots, but fail when tested on websites with different themes, resolutions, layouts, or component libraries. Since the objective of the project is to support functional testing on different web applications, dataset diversity was essential.



\subsection{Second Attempt: Playwright Screenshot Generation}
\label{subsec:playwright_generation}

The second attempt was to generate screenshots using \textbf{Playwright}. This approach made it possible to open web pages automatically, capture screenshots, and create a first collection of realistic interface images.

This method was useful because it allowed the collection of screenshots from real websites and different viewport sizes. It also helped capture desktop and responsive layouts. However, screenshot generation alone did not solve the main problem: the images still needed accurate annotations.

Without reliable labels, the screenshots could not be used directly for YOLO training. This led to the next step: manual and semi-automatic annotation.



\subsection{Third Attempt: Manual Annotation with Roboflow and LabelImg}
\label{subsec:manual_annotation}

Manual annotation was tested using tools such as \textbf{Roboflow}~\cite{roboflow_docs} and \textbf{LabelImg}~\cite{labelimg_repository}. These tools allow bounding boxes to be drawn around UI elements and exported in YOLO format.

Although this approach produced accurate labels for some screenshots, it quickly became difficult to scale. Annotating each button, input field, link, and checkbox manually required a lot of time. This was especially challenging for pages containing many repeated links, menus, forms, and small clickable elements.

Another difficulty was annotation consistency. For example, deciding whether a clickable text should be labeled as a link or a button was not always obvious. Small differences in labeling rules can affect the learning process and reduce model performance.

For these reasons, manual annotation was useful for validation and correction, but it was not sufficient as the main dataset creation strategy.

Figure~\ref{fig:roboflow_manual_labeling} shows one of the manual annotation attempts performed during the dataset preparation phase.

\begin{figure}[H]
	\centering
	\includegraphics[width=0.95\textwidth]{roboflow_manual_labeling.png}
	\caption{Manual annotation attempt using Roboflow}
	\label{fig:roboflow_manual_labeling}
\end{figure}

\subsection{Limitations of Classic Auto-Labeling}
\label{subsec:classic_autolabeling_limits}

Auto-labeling tools were also tested to accelerate the annotation process. However, the results were not immediately reliable. In several cases, the generated boxes were inaccurate, incomplete, or assigned to the wrong class.

This problem is understandable because UI elements are visually diverse. A button may look like a simple text label, an icon, or a traditional rectangle. Input fields may appear with or without borders. Links may be mixed with paragraphs or navigation menus.

The first auto-labeling attempts therefore required significant manual correction. Instead of fully solving the problem, they showed that a more context-aware strategy was necessary.

Figure~\ref{fig:autolabeling_issue} represents an example where automatic labeling did not produce sufficiently reliable annotations.

\begin{figure}[H]
	\centering
	\includegraphics[width=0.95\textwidth]{autolabeling_issue.png}
	\caption{Example of incorrect or incomplete auto-labeling result}
	\label{fig:autolabeling_issue}
\end{figure}

\subsection{Final Solution: DOM-Assisted Auto-Labeling}
\label{subsec:dom_assisted_labeling}

The final solution was to use a DOM-assisted auto-labeling strategy. Instead of relying only on visual auto-labeling, the system used information extracted from the structure of the web page to generate more accurate bounding boxes.

With Playwright, it is possible to inspect DOM elements, retrieve their bounding boxes, and extract useful attributes such as tag name, role, type, placeholder, text content, and accessibility labels. This information was used to identify candidate UI elements and assign them to the appropriate classes.

For example:

\begin{itemize}
	\item \texttt{button} tags and elements with \texttt{role="button"} were labeled as buttons.
	\item \texttt{input} and \texttt{textarea} elements were labeled as input fields.
	\item \texttt{a} tags were labeled as links.
	\item \texttt{select} elements were labeled as dropdowns.
	\item \texttt{input[type="checkbox"]} elements were labeled as checkboxes.
\end{itemize}

This approach significantly improved the quality of the annotations. The labels became more consistent because they were generated from the actual structure of the page instead of being guessed only from pixels. Manual correction was still necessary in some cases, but the amount of manual work was reduced.

Figure~\ref{fig:dom_assisted_labeling} illustrates the DOM-assisted labeling strategy adopted after the first failed attempts. The diagram was created by the author using Mermaid~\cite{mermaid_docs}.

\begin{figure}[H]
	\centering
	\includegraphics[width=0.95\textwidth]{dom_assisted_labeling.png}
	\caption{DOM-assisted labeling approach used to generate UI annotations}
	\label{fig:dom_assisted_labeling}
\end{figure}

\subsection{UI Element Classes}
\label{subsec:ui_element_classes}

The dataset focused on the most common interactive UI elements needed for functional testing. The selected classes were:

\begin{itemize}
	\item \textbf{Button:} Used for actions such as login, submit, search, confirm, or cancel.
	\item \textbf{Input:} Used for text fields, search bars, email fields, password fields, and text areas.
	\item \textbf{Link:} Used for navigation elements and clickable text.
	\item \textbf{Dropdown:} Used for select menus and option lists.
	\item \textbf{Checkbox:} Used for binary selections and form validation.
\end{itemize}

These classes were selected because they cover the most common user interactions in web functional testing.

Table~\ref{tab:ui_classes} describes the UI classes used to annotate the dataset.

\begin{table}[H]
	\centering
	\caption{UI classes used in the vision dataset}
	\label{tab:ui_classes}
	\renewcommand{\arraystretch}{1.35}
	\small
	\begin{tabularx}{\textwidth}{|l|X|X|}
		\hline
		\rowcolor{gray!20}
		\textbf{Class} & \textbf{Description} & \textbf{Examples} \\ \hline
		Button & Clickable component that triggers an action. & Login, Submit, Search \\ \hline
		Input & Field where the user can enter text. & Email, Password, Search field \\ \hline
		Link & Clickable text or navigation element. & Menu link, article link \\ \hline
		Dropdown & Component used to select one option from a list. & Country selector, filter menu \\ \hline
		Checkbox & Binary selection component. & Remember me, Accept terms \\ \hline
	\end{tabularx}
\end{table}

\subsection{Dataset Split and Preparation for Training}
\label{subsec:dataset_split}

After annotation, the dataset was exported in YOLO format. The images and labels were divided into training, validation, and test subsets. This split was necessary to evaluate the model on unseen data and reduce the risk of overfitting.

The training set was used to update the model weights, the validation set was used to monitor performance during training, and the test set was reserved for final evaluation.

Table~\ref{tab:dataset_split} explains the role of each dataset subset used during training and evaluation.

\begin{table}[H]
	\centering
	\caption{Dataset split used for YOLOv8n training}
	\label{tab:dataset_split}
	\renewcommand{\arraystretch}{1.35}
	\small
	\begin{tabularx}{\textwidth}{|l|X|}
		\hline
		\rowcolor{gray!20}
		\textbf{Subset} & \textbf{Role} \\ \hline
		Training set & Used to train the model and update its weights. \\ \hline
		Validation set & Used during training to monitor generalization and tune parameters. \\ \hline
		Test set & Used after training to evaluate the model on unseen examples. \\ \hline
	\end{tabularx}
\end{table}

\section{YOLOv8n Training}
\label{sec:yolov8_training}

After obtaining a more reliable annotated dataset, the next step was to train the object detection model. The objective was to detect UI elements from screenshots and return their class, confidence score, and bounding box coordinates.

\subsection{Why YOLOv8n}
\label{subsec:why_yolov8n}

The selected model was \textbf{YOLOv8n}~\cite{ultralytics_yolov8_docs,yolo_original}. This version was chosen because it offers a good compromise between speed and accuracy. Since the vision module is used during test execution, the model must be fast enough to analyze screenshots without slowing down the whole automation pipeline.

YOLOv8n is also lightweight compared to larger YOLO variants. This makes it more suitable for a prototype that may run in local environments or limited cloud environments.

Table~\ref{tab:yolo_variants} compares the YOLOv8 variants considered during model selection.

\begin{table}[H]
	\centering
	\caption{YOLOv8 variants and project positioning}
	\label{tab:yolo_variants}
	\renewcommand{\arraystretch}{1.35}
	\small
	\begin{tabularx}{\textwidth}{|l|X|X|}
		\hline
		\rowcolor{gray!20}
		\textbf{Model} & \textbf{Advantage} & \textbf{Reason for Selection} \\ \hline
		YOLOv8n & Lightweight and fast. & Selected for real-time screenshot analysis. \\ \hline
		YOLOv8s & More accurate but heavier. & Considered, but less suitable for lightweight execution. \\ \hline
		YOLOv8m/l & Higher capacity models. & Not selected due to resource constraints. \\ \hline
	\end{tabularx}
\end{table}

\subsection{Training Configuration}
\label{subsec:training_configuration}

The training configuration was selected to balance detection quality and computational cost. The model was trained using a standard image size of $640 \times 640$, which is commonly used with YOLO models and preserves enough detail for small UI elements.

Table~\ref{tab:yolo_training_config} presents the final training configuration used for YOLOv8n.

\begin{table}[H]
	\centering
	\caption{YOLOv8n training configuration}
	\label{tab:yolo_training_config}
	\renewcommand{\arraystretch}{1.35}
	\small
	\begin{tabularx}{\textwidth}{|l|l|X|}
		\hline
		\rowcolor{gray!20}
		\textbf{Parameter} & \textbf{Value} & \textbf{Justification} \\ \hline
		Model & YOLOv8n & Lightweight model suitable for real-time detection. \\ \hline
		Image size & 640 & Preserves UI details while keeping inference efficient. \\ \hline
		Epochs & 100 & Allows the model to learn UI element patterns progressively. \\ \hline
		Batch size & 16 & Balanced choice according to available memory. \\ \hline
		Task & Object detection & Required to locate UI elements with bounding boxes. \\ \hline
	\end{tabularx}
\end{table}

\subsection{Loss Function Analysis}
\label{subsec:loss_analysis}

During training, several losses were monitored:

\begin{itemize}
	\item \textbf{Box loss:} Measures the quality of bounding box localization.
	\item \textbf{Class loss:} Measures the quality of class prediction.
	\item \textbf{DFL loss:} Helps refine bounding box boundaries.
\end{itemize}

For this project, box localization is particularly important. If a button is detected but the bounding box is inaccurate, the fallback coordinate click may be placed outside the real element. This would reduce the reliability of the self-healing mechanism.

\subsection{Training Metrics}
\label{subsec:training_metrics}

After training, the model performance was analyzed using the metrics generated by YOLOv8. These metrics include training losses, validation losses, precision, recall, and mean Average Precision.

Figure~\ref{fig:training_results} represents the training metrics obtained after the YOLOv8n training process.

\begin{figure}[H]
	\centering
	\includegraphics[width=1.0\textwidth]{results.png}
	\caption{YOLOv8n training metrics: loss curves, precision, recall, and mAP evolution}
	\label{fig:training_results}
\end{figure}

The curves provide insight into the behavior of the model during training. A decreasing loss indicates that the model is learning useful visual patterns, while stable validation metrics suggest that the model is not only memorizing the training examples.

\subsection{Model Evaluation}
\label{subsec:model_evaluation}

The evaluation focused on whether the model could correctly detect the most important UI components for functional testing. Buttons and input fields were especially important because they are directly used during click and input actions.

The results showed that the model was able to detect several common components correctly. However, some classes remained more difficult, especially when they appeared in small sizes, low contrast, or visually similar styles.

Figure~\ref{fig:confusion_matrix} shows the confusion matrix used to analyze the model behavior across UI classes.

\begin{figure}[H]
	\centering
	\includegraphics[width=0.85\textwidth]{confusion_matrix.png}
	\caption{Confusion matrix of the trained YOLOv8n model}
	\label{fig:confusion_matrix}
\end{figure}

\section{OCR Pipeline}
\label{sec:ocr_pipeline}

Object detection alone is not enough for self-healing execution. YOLO can detect that an element is a button or an input field, but it does not always know which button is the correct one. For example, a page can contain several buttons such as \textit{Login}, \textit{Cancel}, \textit{Search}, and \textit{Submit}.

For this reason, OCR was integrated into the vision pipeline. OCR helps extract visible text from the screenshot and associate it with detected elements.

\subsection{Role of OCR in the Vision Module}
\label{subsec:ocr_role}

The role of OCR is to enrich visual detections with textual information. When the NLP module asks for the \texttt{login button}, the vision module can use OCR to find detected buttons whose visible text is close to \texttt{login}.

This improves the matching process between the semantic target and the visual candidates.

\subsection{OpenCV Preprocessing}
\label{subsec:opencv_preprocessing}

Before applying OCR, screenshots can be preprocessed using \textbf{OpenCV}~\cite{opencv_library}. Preprocessing improves the readability of text and may include grayscale conversion, thresholding, noise reduction, or cropping detected regions.

The objective is not to modify the interface, but to make the text easier for OCR to read.

\subsection{Tesseract OCR}
\label{subsec:tesseract_ocr}

\textbf{Tesseract OCR}~\cite{tesseract_paper,tesseract_docs} was used to extract text from screenshots and detected regions. The extracted text is then cleaned and normalized before being compared with the target extracted from the Gherkin step.

OCR is especially useful for buttons, links, and labels. However, it can be less reliable when the text is small, blurred, partially hidden, or rendered with low contrast.

\subsection{Text Cleaning and Matching}
\label{subsec:text_cleaning_matching}

After OCR extraction, the text is cleaned by removing unnecessary spaces, normalizing case, and filtering irrelevant characters. The cleaned text is then compared with the semantic target using simple matching techniques.

For example, if the target is \texttt{login button}, the matching process prioritizes detections containing the text \texttt{login}. This helps the vision module choose the most relevant element among several candidates.

\section{Vision API Implementation}
\label{sec:vision_api}

The vision module is exposed through API endpoints that allow screenshots to be analyzed and detections to be returned in JSON format.

\subsection{\texttt{/api/ia/analyze-screenshot}}
\label{subsec:api_analyze_screenshot}

The \texttt{/api/ia/analyze-screenshot} endpoint receives a screenshot and applies the vision pipeline. It returns detected UI elements, bounding boxes, confidence scores, and OCR information when available.

This endpoint is useful for testing the vision module independently from the full execution pipeline.

Figure~\ref{fig:analyze_screenshot_request} shows an example request sent to the screenshot analysis endpoint.

\begin{figure}[H]
	\centering
	\includegraphics[width=0.95\textwidth]{analyze_screenshot_request.png}
	\caption{Example request of the \texttt{/api/ia/analyze-screenshot} endpoint}
	\label{fig:analyze_screenshot_request}
\end{figure}

\subsection{Screenshot Input}
\label{subsec:screenshot_input}

The screenshot can come from different sources. It can be uploaded by the user, captured during execution, or generated by Playwright when a selector fails.

In the self-healing process, screenshots are usually captured automatically by the execution service at the moment of failure.

\subsection{Bounding Box Detection}
\label{subsec:bounding_box_detection}

After receiving the screenshot, the YOLOv8n model detects UI elements and returns bounding boxes. Each bounding box contains the position and size of the detected element.

A typical detection includes:

\begin{itemize}
	\item Class name.
	\item Confidence score.
	\item Bounding box coordinates.
	\item Optional OCR text.
\end{itemize}

\begin{lstlisting}[caption={Simplified vision detection response}, label={lst:vision_detection_response}]
	{
		"status": "success",
		"detections": [
		{
			"class_name": "button",
			"confidence": 0.86,
			"bbox": [420, 310, 520, 350],
			"text": "Log In"
		},
		{
			"class_name": "input",
			"confidence": 0.91,
			"bbox": [250, 180, 540, 220],
			"text": "Email"
		}
		]
	}
\end{lstlisting}

\subsection{OCR Enrichment}
\label{subsec:ocr_enrichment}

After detecting bounding boxes, OCR is applied to selected regions in order to extract visible text. This allows the system to associate semantic labels with visual elements.

For example, a detected button with OCR text \texttt{Log In} becomes a stronger candidate when the target is \texttt{login button}.

\subsection{Returned JSON Structure}
\label{subsec:vision_json_structure}

The returned JSON structure is designed to be consumed by the execution service. It provides enough information to select a candidate and perform a coordinate-based action if needed.

The center of the bounding box can be used as the click coordinate:

\[
x_{center} = \frac{x_{min} + x_{max}}{2}
\]

\[
y_{center} = \frac{y_{min} + y_{max}}{2}
\]

This coordinate is then passed to Playwright for fallback execution.

\section{Self-Healing Execution}
\label{sec:self_healing_execution}

The self-healing mechanism is one of the most important parts of \textbf{VisionTest-AI}. Its purpose is to continue execution when the primary selector-based strategy fails.

\subsection{Plan A: Playwright Selector Execution}
\label{subsec:plan_a}

Plan A is the default execution strategy. The system uses selector candidates generated by the semantic mapping module and tries to execute the action with Playwright.

For example, for the step:

\begin{verbatim}
	When I click the login button
\end{verbatim}

the system may try selectors such as:

\begin{itemize}
	\item \texttt{button[name='login']}
	\item \texttt{button[type='submit']}
	\item \texttt{[role='button']:has-text('Log in')}
	\item \texttt{text='login'}
\end{itemize}

If one selector succeeds, the step is considered passed.

\subsection{Plan B: YOLO + OCR + Coordinate-Based Action}
\label{subsec:plan_b}

If Plan A fails, the system activates Plan B. In this case, the execution service captures a screenshot of the current page and sends it to the vision module.

The vision module detects UI elements, extracts OCR text, ranks candidates, and returns the most relevant element. The execution service then uses the center of the selected bounding box to perform a coordinate-based action.

This process allows the system to recover from selector failures in some cases.

Figure~\ref{fig:fallback_example} represents a screenshot captured during a fallback execution, when selector-based execution was not sufficient.

\begin{figure}[H]
	\centering
	\includegraphics[width=0.95\textwidth]{fallback_example.png}
	\caption{Screenshot captured during a fallback execution}
	\label{fig:fallback_example}
\end{figure}

\subsection{Retry Mechanism}
\label{subsec:retry_mechanism}

Web interfaces are often dynamic. Elements may appear after a delay, popups may block interactions, and content may load asynchronously. For this reason, the execution service includes a retry mechanism.

If an action fails, the system can wait briefly, capture a new screenshot, and try again. This improves robustness in pages where the state changes after loading.

\subsection{Popup and Dialog Handling}
\label{subsec:popup_handling}

During testing, popups and dialogs were frequently encountered, especially cookie consent dialogs or login overlays. These elements can block the target UI element and cause execution failure.

To handle this, the system attempts to detect common popup buttons such as \texttt{Accept}, \texttt{Allow}, \texttt{Close}, or \texttt{OK}. When such elements are detected, the system can interact with them before continuing the original test.

\section{Results and Observations}
\label{sec:vision_results_observations}

The vision module was evaluated through both isolated screenshot analysis and full execution scenarios. The objective was not to claim perfect detection, but to verify whether the module could support the self-healing mechanism in realistic cases.

\subsection{Successful Detection Cases}
\label{subsec:successful_detection_cases}

The model produced good results on common UI elements such as buttons and input fields. Detection was especially reliable when the elements had clear borders, sufficient contrast, and standard layouts.

Figure~\ref{fig:vision_detection_result} shows an example of a successful visual detection result.

\begin{figure}[H]
	\centering
	\includegraphics[width=0.95\textwidth]{vision_detection_result.png}
	\caption{Example of successful UI element detection}
	\label{fig:vision_detection_result}
\end{figure}

\subsection{Successful Self-Healing Cases}
\label{subsec:successful_self_healing_cases}

The self-healing mechanism was useful when generated selectors failed but the target element was still visible on the screen. In these cases, the system used screenshot analysis and coordinate-based execution to continue the test.

For example, if the selector for a login button failed because the HTML attributes changed, the vision module could still detect the visible button and provide its coordinates.

Table~\ref{tab:self_healing_success} summarizes examples where the visual fallback helped continue execution after a selector failure.

\begin{table}[H]
	\centering
	\caption{Examples of successful self-healing situations}
	\label{tab:self_healing_success}
	\renewcommand{\arraystretch}{1.35}
	\small
	\begin{tabularx}{\textwidth}{|X|X|X|}
		\hline
		\rowcolor{gray!20}
		\textbf{Situation} & \textbf{Plan A Result} & \textbf{Plan B Result} \\ \hline
		Login button selector not found & Failed selector lookup & Button detected visually and clicked by coordinates \\ \hline
		Search input selector changed & Fill action failed & Input field detected and filled using fallback strategy \\ \hline
		Cookie popup blocked page & Main action blocked & Popup button detected and handled before retry \\ \hline
	\end{tabularx}
\end{table}

\subsection{Failed Cases}
\label{subsec:failed_cases}

Some failures were still observed. The most common causes were OCR errors, low contrast, overlapping elements, and ambiguous targets.

For example, if several buttons had similar text or no visible text, the system could choose the wrong candidate. Similarly, if OCR failed to read the label of a small button, the ranking process became less reliable.

Table~\ref{tab:vision_failed_cases} presents the main failed cases observed during the validation of the vision module.

\begin{table}[H]
	\centering
	\caption{Observed failed cases in the vision module}
	\label{tab:vision_failed_cases}
	\renewcommand{\arraystretch}{1.35}
	\small
	\begin{tabularx}{\textwidth}{|X|X|X|}
		\hline
		\rowcolor{gray!20}
		\textbf{Failure Case} & \textbf{Cause} & \textbf{Possible Improvement} \\ \hline
		Small icon button not detected & Element too small or visually similar to background & Add more icon examples and increase dataset diversity \\ \hline
		Wrong button selected & Several candidates had similar visual/textual features & Improve ranking with DOM context and OCR confidence \\ \hline
		Text not extracted correctly & OCR failed due to low contrast or small font & Improve preprocessing and OCR configuration \\ \hline
		Dropdown not detected reliably & Limited number of dropdown examples in dataset & Add more annotated dropdown samples \\ \hline
	\end{tabularx}
\end{table}

\subsection{Error Analysis}
\label{subsec:vision_error_analysis}

The error analysis shows that the vision module is useful as a fallback mechanism, but it should not be considered a complete replacement for selector-based automation. Playwright selectors remain more precise when the DOM is accessible and stable. The vision module becomes valuable when selectors fail or when the visual state of the page provides stronger information than the HTML structure.

The main improvement needed is dataset diversity. More examples of different UI styles, themes, screen sizes, and component libraries would improve the model's ability to generalize.

\section{Conclusion}
\label{sec:ch4_conclusion}

This chapter presented the visual perception and self-healing execution layer of \textbf{VisionTest-AI}. It first described the dataset preparation process, which evolved from searching for ready-made datasets and manual annotation attempts to a more reliable DOM-assisted auto-labeling approach.

The chapter then presented the YOLOv8n training phase, including the choice of model, training configuration, loss analysis, and evaluation metrics. OCR was also integrated to enrich visual detections with text information and improve target matching.

Finally, the chapter explained the self-healing execution mechanism based on Plan A and Plan B. Plan A relies on Playwright selectors, while Plan B uses YOLO, OCR, and coordinate-based actions when selectors fail.

The results showed that the vision module can improve robustness in some realistic situations, especially when the target element remains visible but selector-based execution fails. However, the module still has limitations related to dataset diversity, OCR reliability, and ambiguous visual targets. The next chapter will focus on the complete system integration, reporting, deployment attempts, and final results of the project.
